% GENERATED FILE. Edit canonical lessons, then run npm run generate:books.
% canonical-source-hash: fnv1a64:d8a2df6b001a9f83
% canonical-lessons: LA-C425-agger, LA-C425-silex, LA-C425-crepido, LA-C425-orbita, LA-C425-sterno

\chapter{Roadbed, Flint, Kerb, Wheel rut, To pave}
\label{ch:la-roadbed-flint-kerb-wheel-rut-to-pave}
\hlchaptermodality{425}

\begin{chapteropening}
\textbf{By the end of this chapter:} \emph{I can describe how a Roman road is built: the raised roadbed, the flint paving, the raised kerb, the wheel ruts, and paving a road.}

Complete the last lesson of chapter 425: I can describe how a Roman road is built: the raised roadbed, the flint paving, the raised kerb, the wheel ruts, and paving a road.
\end{chapteropening}

\section[agger, aggeris]{agger, aggeris — a raised roadbed; a mound, an embankment}

\label{lesson:LA-C425-agger}



\begin{quote}
Before the new one: say the Latin for a winding, then the Latin for to bend.
\end{quote}

\subsection*{You'll want to know: agger, aggeris}
\textbf{agger, aggeris} — \textquotedblleft{}a raised roadbed; a mound, an embankment\textquotedblright{}.

\begin{quote}
Via in aggere altō per campum it.
\end{quote}

\begin{quote}
\emph{The road runs across the plain on a high embankment.}
\end{quote}

Third declension, masculine: \emph{agger, aggeris m.} A Roman road was laid on a raised bank of earth and stones. English \emph{exaggerate} first meant 'to heap up'.

\begin{cousinweb}
\emph{agger} is built on \emph{ad-} 'to' and \emph{gerō} 'to carry' (EtymDB; Lewis \& Short): earth and stones carried to one spot. English \emph{exaggerate} comes from Latin \emph{exaggerāre} 'to heap up', built on it.
\end{cousinweb}

\begin{grammarlens}[title={What you've built}]
One more word for telling, explaining and describing at B1.
\end{grammarlens}

\subsection*{Your turn}
\begin{itemize}
\raggedright
  \item \emph{Say it:} \emph{agger, aggeris}
  \item \emph{Say it:} \emph{agger, aggeris}, once more
  \item \emph{Say it:} a sentence of your own with \emph{agger, aggeris}
  \item \emph{Recall:} say the Latin for a winding, then the Latin for to bend, then say \emph{agger, aggeris} again
\end{itemize}

\begin{tcolorbox}[breakable,colback=teal!4,colframe=teal!35!black,title={Before you move on}]
What does agger, aggeris mean? (\textquotedblleft{}A raised roadbed; a mound, an embankment\textquotedblright{}.) Say it once more.
\end{tcolorbox}
\section[silex, silicis]{silex, silicis — a flint; a hard paving stone}

\label{lesson:LA-C425-silex}



\begin{quote}
Before the new one: say the Latin for to bend, then the Latin for a raised roadbed.
\end{quote}

\subsection*{You'll want to know: silex, silicis}
\textbf{silex, silicis} — \textquotedblleft{}a flint; a hard paving stone\textquotedblright{}.

\begin{quote}
Haec via ē silice facta est, itaque in pluviā quoque bona est.
\end{quote}

\begin{quote}
\emph{This road is made of hard stone, so it is good in the rain too.}
\end{quote}

Third declension, usually masculine: \emph{silex, silicis m.} Roman roads were paved with blocks of it. English \emph{silica} and \emph{silicon} come from it.

\begin{cousinweb}
The origin of \emph{silex} is not settled: Lewis \& Short tentatively link it with a root meaning 'firm', as in \emph{solidus}, and EtymDB gives no parent. EtymDB lists English \emph{silica} and \emph{silicon} among the words made from it.
\end{cousinweb}

\begin{grammarlens}[title={What you've built}]
One more word for telling, explaining and describing at B1.
\end{grammarlens}

\subsection*{Your turn}
\begin{itemize}
\raggedright
  \item \emph{Say it:} \emph{silex, silicis}
  \item \emph{Say it:} \emph{silex, silicis}, once more
  \item \emph{Say it:} a sentence of your own with \emph{silex, silicis}
  \item \emph{Recall:} say the Latin for to bend, then the Latin for a raised roadbed, then say \emph{silex, silicis} again
\end{itemize}

\begin{tcolorbox}[breakable,colback=teal!4,colframe=teal!35!black,title={Before you move on}]
What does silex, silicis mean? (\textquotedblleft{}A flint; a hard paving stone\textquotedblright{}.) Say it once more.
\end{tcolorbox}
\section[crepīdō, crepīdinis]{crepīdō, crepīdinis — a raised edge; a kerb, a raised footway; a quay}

\label{lesson:LA-C425-crepido}



\begin{quote}
Before the new one: say the Latin for a raised roadbed, then the Latin for a flint.
\end{quote}

\subsection*{You'll want to know: crepīdō, crepīdinis}
\textbf{crepīdō, crepīdinis} — \textquotedblleft{}a raised edge; a kerb, a raised footway; a quay\textquotedblright{}.

\begin{quote}
Viātōrēs in crepīdine ambulant, plaustra in mediā viā eunt.
\end{quote}

\begin{quote}
\emph{The travellers walk on the raised footway, the wagons go in the middle of the road.}
\end{quote}

Third declension, feminine: \emph{crepīdō, crepīdinis f.} Any raised edge: the kerb of a street, the base of a building, the stone edge of a quay.

\begin{cousinweb}
\emph{crepīdō} was borrowed from Greek \emph{krēpis}, \emph{krēpidos} 'foundation, base' (EtymDB; Lewis \& Short). The Greek word also meant a kind of boot, the source of Latin \emph{crepida} 'sandal'.
\end{cousinweb}

\begin{grammarlens}[title={What you've built}]
One more word for telling, explaining and describing at B1.
\end{grammarlens}

\subsection*{Your turn}
\begin{itemize}
\raggedright
  \item \emph{Say it:} \emph{crepīdō, crepīdinis}
  \item \emph{Say it:} \emph{crepīdō, crepīdinis}, once more
  \item \emph{Say it:} a sentence of your own with \emph{crepīdō, crepīdinis}
  \item \emph{Recall:} say the Latin for a raised roadbed, then the Latin for a flint, then say \emph{crepīdō, crepīdinis} again
\end{itemize}

\begin{tcolorbox}[breakable,colback=teal!4,colframe=teal!35!black,title={Before you move on}]
What does crepīdō, crepīdinis mean? (\textquotedblleft{}A raised edge; a kerb, a raised footway; a quay\textquotedblright{}.) Say it once more.
\end{tcolorbox}
\section[orbita, orbitae]{orbita, orbitae — a wheel rut, a track}

\label{lesson:LA-C425-orbita}



\begin{quote}
Before the new one: say the Latin for a flint, then the Latin for a raised edge.
\end{quote}

\subsection*{You'll want to know: orbita, orbitae}
\textbf{orbita, orbitae} — \textquotedblleft{}a wheel rut, a track\textquotedblright{}.

\begin{quote}
In viā vetere altae orbitae sunt, et plaustrum lentē it.
\end{quote}

\begin{quote}
\emph{On the old road there are deep ruts, and the wagon goes slowly.}
\end{quote}

First declension, feminine: \emph{orbita, orbitae f.} English \emph{orbit} is the same word: the track a planet runs in.

\begin{cousinweb}
\emph{orbita} is built on \emph{orbis} 'circle, wheel' (EtymDB; Lewis \& Short): the track a wheel leaves. English \emph{orbit} comes from it.
\end{cousinweb}

\begin{grammarlens}[title={What you've built}]
One more word for telling, explaining and describing at B1.
\end{grammarlens}

\subsection*{Your turn}
\begin{itemize}
\raggedright
  \item \emph{Say it:} \emph{orbita, orbitae}
  \item \emph{Say it:} \emph{orbita, orbitae}, once more
  \item \emph{Say it:} a sentence of your own with \emph{orbita, orbitae}
  \item \emph{Recall:} say the Latin for a flint, then the Latin for a raised edge, then say \emph{orbita, orbitae} again
\end{itemize}

\begin{tcolorbox}[breakable,colback=teal!4,colframe=teal!35!black,title={Before you move on}]
What does orbita, orbitae mean? (\textquotedblleft{}A wheel rut, a track\textquotedblright{}.) Say it once more.
\end{tcolorbox}
\section[sternō, sternere]{sternō, sternere — to spread, to lay out; to pave (a road)}

\label{lesson:LA-C425-sterno}



\begin{quote}
Before the new one: say the Latin for a raised edge, then the Latin for a wheel rut.
\end{quote}

\subsection*{You'll want to know: sternō, sternere}
\textbf{sternō, sternere} — \textquotedblleft{}to spread, to lay out; to pave (a road)\textquotedblright{}.

\begin{quote}
Fabrī viam novam saxīs sternunt.
\end{quote}

\begin{quote}
\emph{The workmen pave the new road with stones.}
\end{quote}

\emph{sternō, sternere, strāvī, strātum}. \emph{Viam sternere}, to pave a road; \emph{lectum sternere}, to make a bed. Its participle \emph{strāta} gives English \emph{street}.

\begin{cousinweb}
\emph{sternō} is inherited from Proto-Italic \emph{stornō}, from PIE \emph{sterh\textsubscript{3}-} 'to stretch out, to spread' (EtymDB). Its participle in \emph{via strāta} 'paved road' was borrowed into Germanic and became English \emph{street} (EtymDB).
\end{cousinweb}

\begin{grammarlens}[title={What you've built}]
One more word for telling, explaining and describing at B1.
\end{grammarlens}

\subsection*{Your turn}
\begin{itemize}
\raggedright
  \item \emph{Say it:} \emph{sternō, sternere}
  \item \emph{Say it:} \emph{sternō, sternere}, once more
  \item \emph{Say it:} a sentence of your own with \emph{sternō, sternere}
  \item \emph{Recall:} say the Latin for a raised edge, then the Latin for a wheel rut, then say \emph{sternō, sternere} again
\end{itemize}

\begin{tcolorbox}[breakable,colback=teal!4,colframe=teal!35!black,title={Before you move on}]
What does sternō, sternere mean? (\textquotedblleft{}To spread, to lay out; to pave (a road)\textquotedblright{}.) Say it once more.
\end{tcolorbox}
